% ─────────────────────────────────────────────────────────────────────────────
\section{Background and Related Work}
\label{sec:background}
% ─────────────────────────────────────────────────────────────────────────────

The project sits at the intersection of thermal computer vision, UAV surveillance, and
vision-language reasoning. The literature reviewed during the project showed that these
three areas have each progressed rapidly, but they are not yet naturally aligned for secure,
offline, semantic analysis of thermal imagery.

\subsection{Why Thermal UAV Imagery Is Difficult}

Thermal imagery differs fundamentally from ordinary RGB imagery. It contains heat patterns
instead of color, often has reduced texture, and is especially vulnerable to ambiguity when
small or distant objects are observed from aerial platforms. These properties affect both
detection quality and downstream reasoning quality.

The main challenges are:

\begin{itemize}[nosep]
  \item low contrast and limited fine detail,
  \item absence of color cues and weak texture information,
  \item noise introduced by the sensor and environment,
  \item small target size in wide-area aerial scenes,
  \item domain mismatch between thermal imagery and RGB-pretrained multimodal models.
\end{itemize}

These challenges explain why a naive ``send the whole thermal frame to a VLM'' strategy is
rarely sufficient.

\subsection{Detection-Focused Literature}

Thermal object detection itself is a relatively mature subproblem compared with thermal
semantic reasoning. The YOLO family remains one of the most widely adopted choices for
real-time detection because of its balance between speed and accuracy~\cite{yolov8}. Public
datasets such as FLIR ADAS, KAIST Multispectral, HIT-UAV, DroneVehicle, and
LLVIP~\cite{flir,kaist2015,hituav2022,dronevehicle2022,llvip2021} provide the data basis
for adapting detectors to thermal or cross-modal imagery.

Among these resources, HIT-UAV is particularly relevant because it targets UAV-based
infrared observation and therefore resembles the project scenario more closely than ordinary
ground-camera datasets~\cite{hituav2022}. DroneVehicle is also important because it shows
that aerial cross-modal detection can be treated as an engineering problem with dedicated
datasets and model design choices~\cite{dronevehicle2022}.

The literature review therefore suggested a clear conclusion: \emph{finding targets in thermal
imagery is feasible with specialized detectors}. The harder problem is what happens after the
target is found.

\subsection{Vision-Language Reasoning Models}

Open multimodal models such as Qwen2-VL, LLaVA-NeXT, and InternVL demonstrate
strong reasoning ability on general image understanding tasks~\cite{qwen2vl,llavanext,
internvl2024}. However, their strongest published results come mainly from datasets closer
to RGB photography, document understanding, or general web imagery than to aerial thermal
reconnaissance.

For this reason, the project did not treat the VLM as a replacement for detection. Instead, it
treated the VLM as a second-stage reasoning engine. This design directly addresses three
thermal failure modes that were repeatedly identified during the background study:

\begin{itemize}[nosep]
  \item \textbf{RGB bias:} general multimodal models are optimized for ordinary visual
    concepts rather than heat signatures,
  \item \textbf{small-object blindness:} a full-frame thermal scene contains many distant and
    low-detail targets that are easy for a detector to localize but difficult for a VLM to find,
  \item \textbf{semantic uncertainty:} even when a VLM notices a hot region, it may not
    reliably distinguish whether it represents a vehicle, a person, or irrelevant thermal clutter.
\end{itemize}

This is why the final architecture combined a detection expert and a reasoning expert rather
than pursuing a single-model design.

\subsection{Use of Library and Internet Resources}

Background research for the project relied on a combination of academic and technical
resources:

\begin{itemize}[nosep]
  \item dataset papers and official dataset pages for FLIR ADAS, KAIST, HIT-UAV,
    DroneVehicle, and LLVIP,
  \item official model papers and project pages for YOLOv8, Qwen2-VL, LLaVA-NeXT,
    and InternVL,
  \item official documentation for PyTorch, OpenCV, and NVIDIA Jetson platforms,
  \item professional and ethical references from ACM and IEEE.
\end{itemize}

These resources served different purposes. Academic papers were used to understand the
state of the art and comparable designs. Official documentation was used to understand
component capabilities, deployment constraints, and integration requirements. Ethical codes
and engineering standards were used to frame the system as a decision-support tool rather
than an autonomous authority.

\subsection{Key Takeaway from the Review}

The literature review did not suggest that one new model was needed. Instead, it suggested
that the most practical and defensible solution was \textbf{architectural}: use a detector to
answer ``where is the target?'' and a vision-language model to answer ``what does the target
likely mean?'' This directly motivated the final two-stage ``Gözcü + Analist'' pipeline.
